\chapter{docker-compose.yml}

\begin{lstlisting}[style=PythonStyle, language={}, numbers=left, caption={docker-compose.yml complet}]
version: '3.9'

services:
  db:
    image: postgres:15
    environment:
      POSTGRES_DB: ${DB_NAME:-offres_db}
      POSTGRES_USER: ${DB_USER:-postgres}
      POSTGRES_PASSWORD: ${DB_PASSWORD:-postgres}
    volumes:
      - postgres_data:/var/lib/postgresql/data
    healthcheck:
      test: ["CMD-SHELL", "pg_isready -U ${DB_USER:-postgres}"]
      interval: 10s
      timeout: 5s
      retries: 5

  backend:
    build: ./backend
    command: >
      sh -c "python manage.py migrate --noinput &&
             python manage.py collectstatic --noinput &&
             python manage.py runserver 0.0.0.0:8000"
    environment:
      SECRET_KEY: ${SECRET_KEY}
      DOCUMENT_ENCRYPTION_KEY: ${DOCUMENT_ENCRYPTION_KEY}
    depends_on:
      db:
        condition: service_healthy

  scheduler:
    build: ./backend
    command: >
      sh -c "while true; do
               python manage.py check_deadlines && sleep 3600 || sleep 60;
             done"
    depends_on:
      db:
        condition: service_healthy

  frontend:
    build: ./frontend
    ports:
      - "3000:3000"

  nginx:
    build: ./nginx
    ports:
      - "80:80"
      - "443:443"
    depends_on:
      - backend
      - frontend

volumes:
  postgres_data:
\end{lstlisting}

% -------------------------------------------------------------------

\chapter{Exemple d'entrée AuditLog}

Exemple d'enregistrement dans le journal d'audit pour une opération de publication
d'une offre :

\begin{lstlisting}[style=PythonStyle, language={}, caption={Entrée AuditLog -- publication d'une offre}]
{
  "id": "3fa85f64-5717-4562-b3fc-2c963f66afa6",
  "user": {
    "id": "9b1deb4d-3b7d-4bad-9bdd-2b0d7b3dcb6d",
    "email": "manager@nonedc.ma",
    "full_name": "Karim Benchekroun",
    "role": "MANAGER"
  },
  "action": "UPDATE",
  "resource_type": "OfferCall",
  "resource_id": "7c9e6679-7425-40de-944b-e07fc1f90ae7",
  "changes": {
    "status": "PUBLISHED",
    "published_at": "2025-08-15T09:32:14.000Z"
  },
  "ip_address": "192.168.1.42",
  "timestamp": "2025-08-15T09:32:14.123Z"
}
\end{lstlisting}

% -------------------------------------------------------------------

\chapter{Manuel utilisateur}

\section*{Connexion}

Accéder à l'application via le navigateur. Saisir l'adresse email et le mot de passe.
Cliquer sur \textit{Se connecter}. En cas d'erreur, vérifier les identifiants.

\section*{Créer un appel d'offres}

\begin{enumerate}
    \item Dans le menu latéral, cliquer sur \textbf{Offres} puis \textbf{Nouvelle offre}.
    \item Renseigner le titre, la description, la catégorie, le budget estimatif
    et les dates limites.
    \item Cliquer sur \textbf{Enregistrer}. L'offre est créée en statut BROUILLON.
    \item Pour la publier, ouvrir le détail de l'offre et cliquer sur \textbf{Publier}.
\end{enumerate}

\section*{Évaluer les soumissions}

\begin{enumerate}
    \item Ouvrir le détail de l'offre, onglet \textbf{Soumissions}.
    \item Cliquer sur la soumission à évaluer.
    \item Saisir le score technique et le score financier.
    \item Cliquer sur \textbf{Enregistrer l'évaluation}.
\end{enumerate}

\section*{Portail public (sociétés externes)}

\begin{enumerate}
    \item Accéder à \texttt{https://[domaine]/portal}.
    \item Consulter la liste des appels d'offres ouverts.
    \item Cliquer sur une offre pour voir le descriptif complet.
    \item Cliquer sur \textbf{Soumettre une offre} et remplir le formulaire
    (coordonnées, numéro d'enregistrement, fichier PDF de la proposition).
    \item Conserver le numéro de soumission retourné à titre de référence.
\end{enumerate}

% -------------------------------------------------------------------

\chapter{Résultats du pipeline CI}

% TODO: \insererfigure{images/ci_full.png}{9cm}{Résultats complets du pipeline GitHub Actions}{ci_full}

Le pipeline s'exécute en deux étapes séquentielles :
\begin{enumerate}
    \item \texttt{flake8 .} -- vérification du style Python (0 erreur)
    \item \texttt{pytest --ds=config.settings} -- 12 tests, 0 échec, 0 saut
\end{enumerate}
